\chapter{The sorting bug that rewrote the scoreboard}\label{ch:sortbug}
\begingroup\sloppy % long provenance tags need loose lines rather than overfull ones

\begin{onething}
\code{np.interp} assumes that its sample wavelengths go up and never checks; PRISMA lists its bands going
down. For the implant benchmark's whole early life (2026-10-17 to 18) the evaluator therefore turned every
PRISMA cube it resampled into copies of one or two bands, and flattened every transplanted signature - and the
result everyone believed, ``pretraining is the big step'', was the bug. It was found only when one model was
read two ways on the same real scenes and the two readings were compared.
\end{onething}

This chapter opens the last part of the book, in which the experiments of the project are told as lessons.
Each lesson has the same shape: what was believed, what was tried, what was measured, what it showed, and
the mathematics behind it. The first lesson is the most uncomfortable one, because it is about a result that
looked right, matched what everyone expected, and was wrong for a reason that fits on one line of code.

\section{What was believed}\label{sec:sortbug:believed}

You met the implant benchmark in \cref{ch:benchmarks}: real held-out scenes, called \emph{hosts}, into which
synthetic objects are written with the size, contrast and spectral oddness of the 909 real labelled PRISMA
objects. It was built because the real PRISMA bench (four scenes, 909 objects) could not separate the models:
under residual Mahalanobis (\cref{ch:residualmaha}) every fine-tuned model tied the shipped v0.2.0.

On 2026-10-17 the first version of the implant benchmark, bench-v1 (3,336 objects on 18 hosts\tagM{research/experiments/implant-bench/README.md}),
seemed to give the answer the real bench could not. \Cref{fig:sortbug:stages} shows what it said. On the
easy tier the pAUC (the area under the ROC curve up to 1\% false alarms, divided by 0.01; \cref{ch:measures})
went from 0.257 for the shipped v0.2.0 to 0.342 for F001, the model pretrained on SpectralEarth, and then
crept up to 0.371 through the fine-tuning runs\tagM{research/experiments/indradhanu-finetune/README.md}. The
real bench wandered between 0.148 and 0.159 in no order. The reading written down that day was:
\emph{the largest step is pretraining}.

\coursefig[\linewidth]{iw-56-stages}{pAUC by training stage, as read on 2026-10-17: the implanted bench
(amber, easy and medium tiers) rises from v0.2.0 to FT02c, the real bench (grey) wanders with no order. The
v0.2.0 points are the bug this chapter dissects.}{fig:sortbug:stages}

The second version, bench-v2, built later the same day to fix bench-v1's biases, agreed: v0.2.0 caught 22.1\%
of objects at the top 1\% of scores, F001 29.2\%, FT02c 32.8\%\tagM{research/experiments/implant-bench/README.md}.
Two benchmarks, built differently, telling the same story. And yet one fact did not fit. On the real bench
v0.2.0 tied FT03 (pAUC 0.159 against 0.158), while on bench-v2 it was ten points behind (22.1\% against 32.3\%).
\Cref{ch:benchmarks} left two explanations open: ``a reading artefact'' and ``training overlap''. On 2026-10-18
the reading artefact was confirmed, found to be a two-line bug, and fixed (commit \texttt{b94dbfd}).

\begin{analogy}
A bus timetable is printed with the last bus of the day at the top and the first at the bottom. A passenger
who assumes the usual order checks one thing first: ``is it later than the bottom row?'' The bottom row is the
05:00 bus, so at 09:00, at noon and at 18:00 the answer is always yes - and she is told, every time, ``you
have missed the last bus''. No single answer looks absurd. Only asking at several different times, and getting
the same answer each time, gives it away. Keep this passenger in mind: in \cref{sec:sortbug:contract} you will
see \code{np.interp} behave exactly like her.
\end{analogy}

\section{Linear interpolation, derived}\label{sec:sortbug:interp}

\index{interpolation!linear}
A sensor measures reflectance at a list of wavelengths, its \emph{bands}. To compare sensors, or to move a
spectrum from one sensor's bands to another's, you need values at wavelengths that were not measured.
\emph{Linear interpolation} estimates the value between two neighbouring measurements by the straight line
through them.

\begin{definition}[Linear interpolation]
Given two samples $(x_0, y_0)$ and $(x_1, y_1)$ with $x_0 < x_1$, the linearly interpolated value at a query
point $x$ with $x_0 \le x \le x_1$ is the height at $x$ of the straight line through the two samples.
\end{definition}

Derive the formula. A straight line has a constant slope, so the slope from $(x_0, y_0)$ to $(x, y)$ must equal
the slope from $(x_0, y_0)$ to $(x_1, y_1)$:
\begin{align}
\frac{y - y_0}{x - x_0} &= \frac{y_1 - y_0}{x_1 - x_0}.
\intertext{Multiply both sides by $(x - x_0)$ to isolate $y - y_0$:}
y - y_0 &= \frac{x - x_0}{x_1 - x_0}\,(y_1 - y_0).
\intertext{Name the fraction $t$, the \emph{fraction of the way} from $x_0$ to $x_1$, and add $y_0$ to both sides:}
y &= y_0 + t\,(y_1 - y_0), \qquad t = \frac{x - x_0}{x_1 - x_0}. \label{eq:sortbug:lerp}
\intertext{Expand the bracket and collect $y_0$ to see $y$ as a weighted average of the two samples:}
y &= (1 - t)\,y_0 + t\,y_1. \label{eq:sortbug:lerpweights}
\end{align}
Since $0 \le t \le 1$ inside the interval, the two weights $1-t$ and $t$ are non-negative and add to 1. Outside
the sampled range a rule is needed; the usual one, and NumPy's default, is to hold the end value: below the
smallest sample return its value, above the largest return its value.

\begin{worked}[label={we:sortbug:lerp}]{Interpolating the repo's own test spectrum}
The repository's test (\code{tests/test_evaluation/test_implant.py}, \code{test_descending_band_lists}) uses
samples at 500, 600 and 700\,nm with values 1, 2 and 4. Do these on paper.

At 550\,nm the neighbours are $(500, 1)$ and $(600, 2)$:
\begin{align*}
t &= \frac{550 - 500}{600 - 500} = \frac{50}{100} = 0.5\\
y &= 1 + 0.5 \times (2 - 1)\\
  &= 1 + 0.5\\
  &= 1.5
\end{align*}
At 650\,nm the neighbours are $(600, 2)$ and $(700, 4)$:
\begin{align*}
t &= \frac{650 - 600}{700 - 600} = \frac{50}{100} = 0.5\\
y &= 2 + 0.5 \times (4 - 2)\\
  &= 2 + 1.0\\
  &= 3.0
\end{align*}
At 450\,nm, below the smallest sample, the end value is held: $y = 1$. At 750\,nm, above the largest: $y = 4$.
\end{worked}

\section{The contract of \texttt{np.interp}}\label{sec:sortbug:contract}

\index{np.interp}\index{contract (of a function)}
NumPy's \code{np.interp(x, xp, fp)} takes the query points \code{x}, the sample positions \code{xp} (here
wavelengths) and the sample values \code{fp} (here reflectances). Every function has a \emph{contract}: what it
promises to return, provided its inputs satisfy stated conditions. The conditions are called
\emph{preconditions}. If a precondition is violated, the contract promises nothing.

\begin{definition}[Precondition of \texttt{np.interp}]
The sample positions \code{xp} must be increasing. In the words of NumPy's own docstring (section
``Warnings'', version 2.5.2): ``The x-coordinate sequence is expected to be increasing, but this is not
explicitly enforced. However, if the sequence \code{xp} is non-increasing, interpolation results are
meaningless.''
\end{definition}

Read that carefully. Not an error. Not a warning at run time. \emph{Meaningless} numbers, returned quietly.
Behaviour on decreasing or unsorted \code{xp} is \emph{undefined by contract}: a future NumPy may return
different garbage, and no particular garbage can be relied on. What follows is what NumPy 2.5.2 on the dev
laptop actually returned - an observation, not a promise.

\begin{worked}[label={we:sortbug:descending}]{The same three samples, listed the other way round}
Give \code{np.interp} the identical samples in descending order: \code{xp = (700, 600, 500)},
\code{fp = (4, 2, 1)}. Same physics, same pairs. The values below were produced by
\file{research/textbook/tools/compute_sortbug.py} on NumPy 2.5.2\tagC.

\begin{center}
\begin{tabular}{rcc}
\toprule
query (nm) & correct (\cref{we:sortbug:lerp}) & \texttt{np.interp}, descending \code{xp} \\
\midrule
450 & 1 (held) & 4 \\
500 & 1 & 4 \\
550 & 1.5 & \textbf{1} \\
600 & 2 & \textbf{1} \\
650 & 3.0 & \textbf{1} \\
700 & 4 & \textbf{1} \\
750 & 4 (held) & 1 \\
\bottomrule
\end{tabular}
\end{center}

Every query above 500\,nm got 1, and every query at or below 500\,nm got 4. Only two values ever come out.
\end{worked}

Why only two? The observed table is consistent with an implementation that, for each query $x$, asks two
questions before doing any interpolation, in this order:
\begin{enumerate}
\item Is $x$ beyond the \emph{last} listed sample, $x > \mathtt{xp}[-1]$? If so, return the last listed value
  $\mathtt{fp}[-1]$.
\item Otherwise, is $x$ before the \emph{first} listed sample, $x < \mathtt{xp}[0]$? If so, return the first
  listed value $\mathtt{fp}[0]$.
\item Otherwise, search for the interval that contains $x$ (assuming ascending order) and interpolate.
\end{enumerate}
On an ascending list these checks are the ``hold the end value'' rule. On a descending list they are the
passenger reading the timetable from the bottom.

\begin{keypoint}
On a strictly descending list, under the check order above, step 3 is never reached: every query is answered
by one of the two edge checks, so the output is a step function with only two values.
\end{keypoint}

Prove the key point in two lines. Let the list be descending, so its last entry is its smallest and its first
is its largest: $\mathtt{xp}[-1] < \mathtt{xp}[0]$. Take any query $x$.
\begin{itemize}
\item If $x > \mathtt{xp}[-1]$, check 1 fires and returns $\mathtt{fp}[-1]$.
\item If not, then $x \le \mathtt{xp}[-1]$. Chain the inequality with $\mathtt{xp}[-1] < \mathtt{xp}[0]$ to
  get $x < \mathtt{xp}[0]$, so check 2 fires and returns $\mathtt{fp}[0]$.
\end{itemize}
Every $x$ falls in one of the two cases, so step 3 never runs. With $\mathtt{xp} = (700, 600, 500)$: every
$x > 500$ gets $\mathtt{fp}[-1] = 1$ and every $x \le 500$ gets $\mathtt{fp}[0] = 4$ - exactly the table.
The proof also covers a list made of several descending runs, as long as its last entry is smaller than its
first; you will need that for the PRISMA hosts in \cref{sec:sortbug:prisma}.

\coursefig[\linewidth]{ad-22-interp-order}{The same three samples (500, 600, 700\,nm; values 1, 2, 4). (a) Listed
ascending, \texttt{np.interp} draws the line: 1.5 at 550\,nm and 3.0 at 650\,nm. (b) Listed descending, it
returns the last listed value, 1, for every wavelength above 500\,nm: the spectrum goes flat.}{fig:sortbug:interp}

\section{Resampling a cube is a matrix}\label{sec:sortbug:matrix}

\index{resampling!as a matrix}
The evaluator's \code{resample} in \code{app/evaluation/datasets.py} moves a whole cube of shape
\shape{(C, H, W)} (C bands, H rows, W columns) onto a fixed grid of G wavelengths. It does not call
\code{np.interp} once per pixel. It builds one \emph{weight matrix} $\mat{M}$ of shape \shape{(G, C)} and
multiplies every pixel's spectrum by it. Why is that allowed?

\paragraph{Interpolation is linear in the values.}
Fix the sample positions and the query point. Equation \eqref{eq:sortbug:lerpweights} says the output is
$(1-t)\,y_0 + t\,y_1$ - a weighted sum of the sample values, with weights that depend only on positions. So if
you interpolate two value lists $\vect{f}$ and $\vect{g}$ and add the results, you get the same as interpolating
$\vect{f} + \vect{g}$; and scaling the values by a number scales the output by that number. A map with those
two properties is \emph{linear}, and every linear map from $C$ numbers to $G$ numbers is a $G \times C$ matrix.

\paragraph{Finding the matrix column by column.}
Write a pixel's spectrum as a sum of \emph{unit vectors}: $\vect{e}_b$ is the list with 1 in position $b$ and 0
elsewhere, so a spectrum $\vect{y} = (y_1, \ldots, y_C)$ equals $\sum_b y_b\,\vect{e}_b$. By linearity,
\begin{align}
\operatorname{interp}(\vect{y}) &= \operatorname{interp}\Big(\sum_b y_b\,\vect{e}_b\Big)
  = \sum_b y_b\,\operatorname{interp}(\vect{e}_b).
\intertext{The right-hand side is a matrix-vector product whose $b$-th column is $\operatorname{interp}(\vect{e}_b)$:}
\operatorname{interp}(\vect{y}) &= \mat{M}\vect{y}, \qquad \mat{M}_{:,b} = \operatorname{interp}(\text{grid};\ \mathtt{wl},\ \vect{e}_b).
\end{align}
That is the code, line for line:

\begin{lstlisting}
# app/evaluation/datasets.py, before b94dbfd
# M: (G, C); column i is the interpolation of unit vector e_i
M = np.stack([np.interp(grid, wl, np.eye(len(wl))[i]) for i in range(len(wl))], 1)
return np.einsum("gb,bhw->ghw", M, cube, optimize=True)   # (C, H, W) -> (G, H, W)
\end{lstlisting}

\code{np.eye(len(wl))[i]} is $\vect{e}_i$; \code{np.stack(..., 1)} puts the results side by side as columns;
\code{einsum("gb,bhw->ghw")} multiplies every pixel's $C$-vector by $\mat{M}$.

\begin{worked}[label={we:sortbug:row}]{One row of the resample matrix, right and wrong}
Three scene bands listed 700, 600, 500\,nm (descending, as PRISMA lists them); one grid wavelength, 550\,nm.
The pixel reads $(4, 2, 1)$ in that listed order.

The correct row puts weight on the two bands that bracket 550\,nm. From \cref{we:sortbug:lerp}, $t = 0.5$ of
the way from 500 to 600, so 600\,nm gets $t = 0.5$ and 500\,nm gets $1 - t = 0.5$:
\begin{center}
\begin{tabular}{lccc}
\toprule
band as listed & 700 & 600 & 500 \\
\midrule
correct row for 550\,nm & 0.0 & 0.5 & 0.5 \\
unsorted row for 550\,nm & 0.0 & 0.0 & 1.0 \\
pixel value & 4 & 2 & 1 \\
\bottomrule
\end{tabular}
\end{center}
Where does the unsorted row come from? Column $b$ is $\operatorname{interp}(550;\ (700,600,500),\ \vect{e}_b)$.
By the key point of \cref{sec:sortbug:contract}, $550 > 500$ (the last listed), so each call returns the last
entry of $\vect{e}_b$: 0 for $b = 700$, 0 for $b = 600$, 1 for $b = 500$.

Apply each row to the pixel:
\begin{align*}
\text{correct:}\quad & 0 \times 4 + 0.5 \times 2 + 0.5 \times 1\\
 &= 0 + 1.0 + 0.5\\
 &= 1.5\\
\text{unsorted:}\quad & 0 \times 4 + 0 \times 2 + 1 \times 1\\
 &= 1
\end{align*}
\end{worked}

\paragraph{The fix: sort the pairs, keep them together.}
\index{argsort}
The fix sorts the wavelengths and builds the columns on the sorted list, then writes each column back to the
position of the band it belongs to:

\begin{lstlisting}
# app/evaluation/datasets.py, since b94dbfd
order = np.argsort(wl)                # (C,) band positions in ascending order
M = np.zeros((len(grid), len(wl)))    # (G, C)
M[:, order] = np.stack([np.interp(grid, wl[order], np.eye(len(wl))[i])
                        for i in range(len(wl))], 1)
\end{lstlisting}

\code{np.argsort(wl)} returns the \emph{permutation} that sorts \code{wl}: the list of positions such that
\code{wl[order]} is ascending. Sorting does not change the data. Each wavelength keeps its own value; only the
order in which the pairs are listed changes. A function sampled at a set of points is the same function however
you list the points - the bug was never in the data, only in a precondition left unmet.

\begin{worked}[label={we:sortbug:fix}]{The sorted resample, by hand}
Bands listed $\mathtt{wl} = (700, 600, 500)$; grid $(550, 650)$; pixel $(4, 2, 1)$ in listed order.
\begin{enumerate}
\item Sort. The smallest wavelength, 500, sits at position 2; then 600 at position 1; then 700 at position 0. So
  $\mathtt{order} = (2, 1, 0)$ and $\mathtt{wl[order]} = (500, 600, 700)$, ascending.
\item Build columns on the sorted list. Sorted column 0 (500\,nm): weight $0.5$ at 550, $0$ at 650. Sorted
  column 1 (600\,nm): $0.5$ at 550, $0.5$ at 650. Sorted column 2 (700\,nm): $0$ at 550, $0.5$ at 650.
\item Write back: sorted column $i$ goes to original column $\mathtt{order}[i]$. Sorted 0 $\to$ original 2,
  sorted 1 $\to$ original 1, sorted 2 $\to$ original 0:
  \[
  \mat{M} = \begin{pmatrix} 0 & 0.5 & 0.5 \\ 0.5 & 0.5 & 0 \end{pmatrix}
  \quad\text{(columns: 700, 600, 500 nm).}
  \]
\item Multiply. Row 1: $0 \times 4 + 0.5 \times 2 + 0.5 \times 1 = 0 + 1 + 0.5 = 1.5$. Row 2:
  $0.5 \times 4 + 0.5 \times 2 + 0 \times 1 = 2 + 1 + 0 = 3.0$.
\end{enumerate}
The old code gives $\mat{M} = \begin{psmallmatrix} 0 & 0 & 1 \\ 0 & 0 & 1\end{psmallmatrix}$ and the
output $(1, 1)$. Both matrices were confirmed by \file{research/textbook/tools/compute_sortbug.py}\tagC.
\end{worked}

\section{What it did to real PRISMA band lists}\label{sec:sortbug:prisma}

PRISMA has two detectors, a visible and near-infrared (VNIR) one and a short-wave infrared (SWIR) one, and its
files list the bands of each detector from long to short wavelength. Two real lists met the bug.

\paragraph{The bench objects' list.} The 185 bands on which the real PRISMA bench objects were measured
(\file{C:/es/implant/real_v2.json}, field \texttt{wavelengths\_nm}) start at 2477.05\,nm and end at
453.39\,nm\tagM{C:/es/implant/real_v2.json}. The 165-band common grid of \cref{ch:sensors} runs from 460 to
2450\,nm in 10\,nm steps. Every grid wavelength is longer than 453.39\,nm, the last band listed. By the key
point, every row of the old $\mat{M}$ has a single 1, in column 184 - the 453\,nm band. Running the old line
on the real list confirms it: only column 184 is ever used\tagC.

\begin{keypoint}
The fixed-grid model was handed 165 copies of the same blue image: no spectrum to reconstruct, only one band's
spatial texture.
\end{keypoint}

\paragraph{The implant hosts' lists.} A PRISMA host's \file{wavelengths.npy} (for example
\file{C:/es/implant/bench-v2.1/prisma/PRS_L2D_STD_20230311055340/}) holds 234 bands in two descending runs: the
VNIR detector from 977.37 down to 406.99\,nm, then the SWIR detector from 2497.12 down to
943.36\,nm\tagM{C:/es/implant/bench-v2.1/prisma/PRS_L2D_STD_20230311055340/wavelengths.npy}. The list is not
descending overall, but its last entry (943.36) is smaller than its first (977.37), which is all the proof in
\cref{sec:sortbug:contract} needs:
\begin{align*}
\text{grid wavelength} > 943.36\,\text{nm} &\;\to\; \text{the 943\,nm band (last listed)},\\
\text{grid wavelength} \le 943.36\,\text{nm} &\;\to\; \text{the 977\,nm band (first listed)}.
\end{align*}
Count them on the common grid. The grid skips the water-vapour window from 912 to 978\,nm
(\code{build_common_wavelength_grid} in \code{app/models/dataset/vendables.py}), so its wavelengths at or below
943.36\,nm are 460, 470, \ldots, 910:
\[
\frac{910 - 460}{10} + 1 = \frac{450}{10} + 1 = 45 + 1 = 46 \text{ grid bands.}
\]
So 46 grid bands became copies of the 977\,nm image and the other $165 - 46 = 119$ copies of the 943\,nm image
(the script agrees: 46, and only columns 0 and 233 are used\tagC). Every grid band was one of two near-infrared
pictures.

\begin{sourcenote}[: which models the reading bug touched]
Only fixed-grid models go through \code{resample}: the shipped v0.2.0 reads the 165-band grid. The
wavelength-reading models (F001 and every FT run) read the host's own bands directly and never touched it. And
only descending lists trip it: the EnMAP and AVIRIS-NG hosts were read correctly. The real PRISMA bench was also
safe: for fixed-grid models the evaluator reloads its scenes through the product's \code{band_filter_apply},
whose resampler (\code{app/utils/data_transformations/spectral_resampler.py}, line 70) begins
\texttt{\# Sort source wavelengths ascending}. The product always sorted; the evaluator re-implemented a faster
resample and left the sort out.
\end{sourcenote}

\section{What it did to the signature implants}\label{sec:sortbug:signature}

\index{implant!signature}
A \emph{signature} implant (\cref{ch:benchmarks}) takes a real object's difference spectrum $\vect{d} =
\vect{o} - \vect{r}$ - the object's mean spectrum minus the mean of a ring of pixels around it - measured on
the 185 bench bands, and carries it to the host's bands by interpolation. Before \texttt{b94dbfd} that was a bare
\code{np.interp(wl, rw, d)}, in the builder's \code{place_and_implant} and in \code{implant_tile} alike, with
\code{rw} the descending bench list. Same mechanism, same damage.

\begin{worked}[label={we:sortbug:signature}]{A signature flattened}
A difference spectrum listed descending: $+0.06$ at 700\,nm, $+0.03$ at 600\,nm, $-0.02$ at 500\,nm. The host
has bands at 520, 580 and 650\,nm. Do it on paper.

Correct, by \eqref{eq:sortbug:lerp}:
\begin{align*}
520\,\text{nm}:\quad t &= \tfrac{520 - 500}{600 - 500} = 0.2, & y &= -0.02 + 0.2 \times (0.03 - (-0.02))\\
 & & &= -0.02 + 0.01 = -0.01\\
580\,\text{nm}:\quad t &= \tfrac{580 - 500}{600 - 500} = 0.8, & y &= -0.02 + 0.8 \times 0.05\\
 & & &= -0.02 + 0.04 = 0.02\\
650\,\text{nm}:\quad t &= \tfrac{650 - 600}{700 - 600} = 0.5, & y &= 0.03 + 0.5 \times (0.06 - 0.03)\\
 & & &= 0.03 + 0.015 = 0.045
\end{align*}
Unsorted: all three host bands are above 500\,nm (the last listed), so each gets the last listed value,
$-0.02$. The real shape - dark in the blue, bright towards the red - arrives as a flat step $(-0.02, -0.02,
-0.02)$. Both results match the script\tagC.
\end{worked}

For the real bench list, every host band longer than 453.39\,nm received $\vect{d}$ at 453\,nm, one number;
the few host bands at or below 453\,nm received $\vect{d}$ at 2477\,nm. bench-v2's later steps
(\code{match_foreignness} and the strength solve of \cref{ch:benchmarks}) then rescaled that step, so the
garbled implants still hit their target contrast and foreignness - but not the shape of any real object.

\begin{keypoint}
The signature bug hit every signature implant in bench-v1 and bench-v2, for every model, and half of the
implants FT05 trained on (\cref{ch:ftruns}).
\end{keypoint}

The fix is a three-line function:
\begin{lstlisting}
# app/evaluation/implant.py
def interp_spectrum(wl, src_wl, values):
    src_wl, values = np.asarray(src_wl, float), np.asarray(values, float)
    order = np.argsort(src_wl)                     # ascending, pairs kept together
    return np.interp(wl, src_wl[order], values[order])
\end{lstlisting}

\section{How it hid in plain sight}\label{sec:sortbug:hid}

Five things kept the bug invisible for two days. Each is a habit worth breaking.
\begin{enumerate}
\item \textbf{No error.} \code{np.interp} returned numbers of the right shape and range. Nothing crashed,
  nothing was NaN. A copy of a real band is a perfectly plausible-looking image.
\item \textbf{A warning that sounded like an excuse.} The resample's docstring said implant hosts are read
  ``close to, not identical with'' the product. A difference was expected, so a difference was not alarming.
\item \textbf{A story that fitted.} On bench-v1 the models lined up in training order with v0.2.0 at the bottom,
  and ``pretraining helps'' is what everyone expected. The Indradhanu course even noted that the model and the
  band reading changed together between v0.2.0 and F001 - and read the step as the model's.
\item \textbf{A cross-check that was run and misread.} The implant-bench README now says: ``An earlier note here
  that the two readings were identical was a misread.'' The saved check said 15.8\% against 44.2\%.
\item \textbf{The tests used the wrong orderings.} Every test of the new resample used ascending lists, where
  sorted and unsorted code agree, until \code{test_descending_band_lists} was added with the fix. It feeds the
  same cube with the wavelengths reversed (\code{up[::-1]}) and asserts the result equals the ascending one,
  and that \code{interp_spectrum} returns 1.5 and 3.0 at 550 and 650\,nm from a descending list.
\end{enumerate}

\section{The cross-check that found it}\label{sec:sortbug:crosscheck}

\index{cross-check}
For a fixed-grid model the real PRISMA bench can be read two ways: through the product's
\code{band_filter_apply} path (what the evaluator does for PRISMA bench scenes), or through the generic
\code{resample} (what implant hosts get). If the evaluator is right, the two must agree on the same scenes. A
one-off script, \file{C:/es/eval/resample_check.py}, swapped the reader and scored v0.2.0
(\cref{tab:sortbug:crosscheck}).

\begin{table}[htbp]
\centering
\caption{v0.2.0 on the real PRISMA bench (909 objects), residual Mahalanobis, read three ways.}
\label{tab:sortbug:crosscheck}
\small
\begin{tabular}{lccccc}
\toprule
reading path & caught top 1\% & top 0.1\% & pAUC to 1\% & ROC-AUC & TPR at 0.1\% \\
\midrule
product path & 44.2\% & 14.1\% & 0.159 & 0.854 & 0.058 \\
generic resample, before fix & 15.8\% & 5.6\% & 0.094 & 0.752 & 0.023 \\
generic resample, after fix & 44.2\% & 14.1\% & 0.159 & 0.854 & 0.058 \\
\bottomrule
\end{tabular}

\smallskip
{\footnotesize Sources: \tagM{C:/es/eval/resample_check/report.md}; \tagM{research/experiments/implant-bench/README.md} (Corrections, 2026-10-18).}
\end{table}

After the fix the two paths give the same five numbers. In objects, count what the bug had cost v0.2.0.

\begin{worked}[label={we:sortbug:objects}]{Objects the bug hid}
At the top 1\%:
\begin{align*}
0.442 - 0.158 &= 0.284\\
0.284 \times 909 &= 258.2 \approx 258 \text{ objects}
\end{align*}
At the top 0.1\%:
\begin{align*}
0.141 - 0.056 &= 0.085\\
0.085 \times 909 &= 77.3 \approx 77 \text{ objects}
\end{align*}
(Script check: 258.156 and 77.265\tagC.)
\end{worked}

\begin{keypoint}
When a benchmark has its own data path, run one model through both paths on data where both exist, and treat any
gap as a bug until it is explained.
\end{keypoint}

\section{The scoreboard, read correctly: bench-v2.1}\label{sec:sortbug:board}

bench-v2 was rebuilt with the fix as \emph{bench-v2.1} (\texttt{--version 2 --psf-sigma 1.4 --foreign-tol 1.2
--seed 0}): 18 hosts, 2,817 objects (signature 985, material 927, host
905)\tagM{research/experiments/implant-bench/README.md}. The foreignness tolerance was also tightened from 1.5 to
1.2, the candidate fix for bench-v2's marginal acceptance miss. \Cref{tab:sortbug:v21} is the board.

\begin{table}[htbp]
\centering
\caption{bench-v2.1, residual Mahalanobis. ``Caught'' is the share of the 2,817 objects with a pixel within
one step of a top-scoring pixel (\cref{ch:measures}).}
\label{tab:sortbug:v21}
\small\setlength{\tabcolsep}{3.5pt}
\begin{tabular}{lcccccc}
\toprule
model & top 1\% & top 0.1\% & pAUC to 1\% & TPR at 0.1\% & EnMAP hosts & PRISMA hosts \\
\midrule
\textit{chance} & 10.5\% & 1.1\% & & & & \\
v0.2.0 (read correctly) & \textbf{35.5\%} & 12.6\% & 0.168 & 0.082 & 21.6\% & \textbf{42.2\%} \\
F001 ep109 & 33.1\% & 11.7\% & 0.157 & 0.078 & 20.5\% & 37.7\% \\
FT02c ep10 & 35.1\% & 12.8\% & 0.169 & 0.086 & 20.7\% & 40.3\% \\
FT03 ep10 & 34.9\% & 12.7\% & 0.168 & 0.086 & 20.6\% & 40.1\% \\
FT04 ep10 & 35.1\% & 12.7\% & 0.169 & 0.083 & 20.3\% & 40.4\% \\
FT05 ep5 & 35.4\% & 12.8\% & 0.170 & 0.085 & 20.8\% & 40.7\% \\
PCA-RX (no model) & 33.2\% & 9.8\% & 0.154 & 0.069 & 19.8\% & 37.7\% \\
\bottomrule
\end{tabular}

\smallskip
{\footnotesize Sources: \tagM{research/experiments/implant-bench/README.md} (bench-v2.1); FT04's row
\tagM{research/experiments/indradhanu-finetune/FT04-final-joint/README.md} and \tagM{C:/es/eval/ft04/report.md}.}
\end{table}

\begin{drift}
\Cref{tab:sortbug:v21} and \cref{fig:sortbug:bars} are the 2026-10-18 board (FT04 added 2026-10-19). Later runs
(FT06 to FT09, Indradhanu Heavy; \cref{ch:uncertainty,ch:decoderwork,ch:heavy}) are scored on the same bench and
are not shown here. Any rebuild of the bench moves every row.
\end{drift}

Read it against the old boards (\cref{fig:sortbug:bars}). v0.2.0 caught 19.0\% on bench-v1, 22.1\% on bench-v2,
35.5\% on bench-v2.1. FT03 caught 26.9\%, 32.3\%, 34.9\%\tagM{research/experiments/implant-bench/README.md}.

\begin{figure}[htbp]
\centering
\begin{tikzpicture}
\begin{axis}[
  ybar, bar width=12pt, width=0.92\linewidth, height=6.5cm,
  ymin=0, ymax=42, ylabel={caught at top 1\% (\%)},
  symbolic x coords={bench-v1,bench-v2,bench-v2.1}, xtick=data,
  enlarge x limits=0.22, legend style={at={(0.5,1.03)},anchor=south,legend columns=4,font=\small},
  nodes near coords, nodes near coords style={font=\tiny, rotate=90, anchor=west},
  every axis plot/.append style={draw=black!60}]
\addplot[fill=loss!70] coordinates {(bench-v1,19.0) (bench-v2,22.1) (bench-v2.1,35.5)};
\addplot[fill=spec!60] coordinates {(bench-v1,25.2) (bench-v2,29.2) (bench-v2.1,33.1)};
\addplot[fill=spat!50] coordinates {(bench-v1,26.9) (bench-v2,32.3) (bench-v2.1,34.9)};
\addplot[fill=maskgrey!60] coordinates {(bench-v1,30.2) (bench-v2,30.2) (bench-v2.1,33.2)};
\legend{v0.2.0 (fixed grid), F001, FT03, PCA-RX}
\end{axis}
\end{tikzpicture}
\caption{Objects caught at the top 1\% on the three implant benches, drawn to scale. The wavelength-reading models
(F001, FT03) and PCA-RX move a few points with each rebuild; v0.2.0, whose reading was the only one fixed, jumps
13.4 points from bench-v2 to bench-v2.1. Its bench-v1 and bench-v2 bars measured a broken resample.}
\label{fig:sortbug:bars}
\end{figure}

\begin{worked}[label={we:sortbug:split}]{How much of v0.2.0's jump is the fix?}
v0.2.0 moved
\[
35.5 - 22.1 = 13.4 \text{ points}
\]
from bench-v2 to bench-v2.1. FT03, which never went through the broken resample, moved
\[
34.9 - 32.3 = 2.6 \text{ points}
\]
with the rebuild alone (correct signatures, tighter tolerance). If the rebuild gave v0.2.0 the same 2.6 points,
the fix accounts for
\[
13.4 - 2.6 = 10.8 \text{ points}.
\]
This split is rough: the rebuild changed the bench for every model, but not by exactly the same amount.
\end{worked}

\paragraph{How big is noise? A rough ruler.}
\index{standard error!binomial}
Is v0.2.0's 35.5\% really ahead of FT03's 34.9\%? Treat each object as a coin that a model catches with some
probability $p$. Write $X_i = 1$ if object $i$ is caught and 0 if not; this is a \emph{Bernoulli} variable.
Derive its mean and variance:
\begin{align}
\E[X_i] &= 1 \cdot p + 0 \cdot (1-p) = p, \notag\\
\E[X_i^2] &= 1^2 \cdot p + 0^2 \cdot (1-p) = p, \notag\\
\operatorname{Var}[X_i] &= \E[X_i^2] - (\E[X_i])^2 = p - p^2 = p(1-p). \notag
\intertext{The share caught is the average $\bar X = \frac{1}{n}\sum_i X_i$. For independent objects the variance
of a sum is the sum of variances, and dividing by $n$ divides the variance by $n^2$:}
\operatorname{Var}[\bar X] &= \frac{1}{n^2} \cdot n\,p(1-p) = \frac{p(1-p)}{n}, \qquad
\operatorname{SE} = \sqrt{\frac{p(1-p)}{n}}. \label{eq:sortbug:se}
\end{align}

\begin{worked}[label={we:sortbug:se}]{Is 35.5\% against 34.9\% a difference?}
With $p = 0.35$ and $n = 2{,}817$:
\begin{align*}
p(1-p) &= 0.35 \times 0.65 = 0.2275\\
0.2275 / 2{,}817 &= 0.0000808\\
\sqrt{0.0000808} &= 0.0090 \quad\text{(about 0.9 points)}\\
\text{gap} &= 35.5 - 34.9 = 0.6 \text{ points} = 0.6 / 0.9 = 0.7 \text{ standard errors}
\end{align*}
Under one standard error: a tie. (Script: $\operatorname{SE} = 0.00899$, gap 0.67\,SE\tagC.)
\end{worked}

Two models scored on the \emph{same} objects are correlated, so the honest uncertainty of their difference is
smaller than this single-model ruler suggests. \Cref{ch:bootstrap} builds the proper paired object bootstrap; use
\eqref{eq:sortbug:se} only as a first, cautious check.

Three more readings of \cref{tab:sortbug:v21}:
\begin{itemize}
\item \textbf{Every model stays above PCA-RX at the top 0.1\%} (11.7--12.8\% against 9.8\%), as on the real
  bench - the ordering test of \cref{ch:benchmarks} still passes. At the top 1\% F001 and PCA-RX tie (33.1\%
  against 33.2\%).
\item \textbf{EnMAP hosts sit near 21\% for every model and for PCA-RX.} Unexplained on bench-v1, v2 and v2.1.
\item \textbf{The easy tier is too foreign.} Per-tier foreignness on PRISMA hosts is 1.05$\times$ / 1.15$\times$ /
  1.60$\times$ real (subtle / medium / easy); halo 0.92--1.06$\times$\tagM{research/experiments/implant-bench/README.md}.
  The easy tier is outside the 25\% acceptance limit.
\end{itemize}

\begin{worked}[label={we:sortbug:selection}]{The selection effect behind the easy tier}
This is an illustration, not bench data\tagA. Four easy targets want foreignness 0.001, 0.002, 0.004 and 0.008.
With the tolerance at 1.2, suppose the realised object can only get within 1.2$\times$ of the two higher targets, so
the two lower spots are skipped. The median of four numbers is the mean of the middle two:
\begin{align*}
\text{targets:}\quad & 0.001,\ 0.002,\ 0.004,\ 0.008 \qquad \text{median} = (0.002 + 0.004)/2 = 0.003\\
\text{survivors:}\quad & 0.004,\ 0.008 \qquad \text{median} = (0.004 + 0.008)/2 = 0.006
\end{align*}
Nothing was implanted wrongly; the population that remains is not the one aimed at.
\end{worked}

\begin{keypoint}
Read correctly, v0.2.0 ties or beats every new model on bench-v2.1, as it ties them on the real bench: the
``pretraining is the big step'' result (v0.2.0 19--22\% to F001 25--29\%) was the reading bug.
\end{keypoint}

\section{Lessons for trusting a benchmark}\label{sec:sortbug:lessons}

\Cref{ch:benchmarks}'s ordering test asked: does the benchmark reproduce something the real labels already told
you? It was applied to PCA-RX, and passed. It was not applied to v0.2.0 - and v0.2.0 failed it on both bench-v1
and bench-v2 (a tie on real objects, far behind on implants). That failure was the bug announcing itself. So,
before you trust a benchmark's ranking of models:
\begin{enumerate}
\item \textbf{Check every model's position against the real labels, not only the baseline's.} One model that
  the real bench and the benchmark disagree about is a flag on the benchmark's data path.
\item \textbf{Cross-check any re-implemented data path against the product's}, on data where both exist
  (\cref{sec:sortbug:crosscheck}).
\item \textbf{Test with real input orderings.} PRISMA's descending list is a property of real files; a test
  written with \code{[500, 600, 700]} cannot see it.
\item \textbf{Know your preconditions.} If a library function's contract says ``meaningless'' for some input,
  either enforce the precondition (sort) or assert it (\code{np.all(np.diff(xp) > 0)}).
\item \textbf{Give a result you hoped for the scrutiny you would give one you feared.}
\end{enumerate}

The same descending list surfaced once more, in the product: when Indradhanu V2 began reading native PRISMA bands
(\cref{ch:shipping}), \code{scene_segmentation}'s brightness - ``the mean of band 0 up to the VNIR edge'' - turned
out to assume ascending order. On native PRISMA, band 0 is the far SWIR. That fix is part of the next two
chapters' story.

\begin{confusions}
\confusion{The bug made every model's implant numbers wrong.}{Two different effects. The \emph{reading} bug hit
fixed-grid models (v0.2.0) on PRISMA hosts only. The \emph{signature} bug garbled half (bench-v1) or a third
(bench-v2) of the implants for every model alike. The wavelength-reading models' rows moved with the rebuild, not
with the reading fix.}
\confusion{The real bench was wrong too.}{No. For fixed-grid models the evaluator reloads PRISMA bench scenes through
the product's \texttt{band\_filter\_apply}, which always sorted. v0.2.0's real-bench 0.159 never changed.}
\confusion{Sorting the wavelengths reorders the data.}{It reorders the pairs. Band 453\,nm keeps its own values; it
is only listed first instead of last. The function being sampled is unchanged.}
\confusion{NumPy guarantees the two-value step on descending input.}{It guarantees nothing: the contract calls the
result meaningless. The step is what NumPy 2.5.2 returned on the dev laptop, consistent with the check order of
\cref{sec:sortbug:contract}. Do not rely on it, and do not write a test that expects it.}
\confusion{bench-v2.1 says v0.2.0 is the best model.}{It says v0.2.0 ties the best models: 0.6 points ahead of FT03
is under one standard error. No new model measurably beats it on detection.}
\confusion{bench-v2.1 failed acceptance, so it cannot be used.}{The easy tier's foreignness misses the 25\% limit,
by a known selection effect. The ordering test still holds. Use it as a ranking with that caveat stated, and keep
the real bench as the judge.}
\end{confusions}

\begin{checkyourself}
\item\label{ex:sortbug:1} By hand: samples at 800, 700, 600\,nm with values 0.30, 0.20, 0.10. What is the correct
  value at 650\,nm, and what does \texttt{np.interp} return (NumPy 2.5.2, the check order of
  \cref{sec:sortbug:contract}) if the samples are listed 800, 700, 600? What does it return at 590\,nm?
\item\label{ex:sortbug:2} Bands listed 700, 600, 500\,nm. Write the correct resample row for a grid wavelength of
  620\,nm, then the row the old unsorted code built, and apply both to a pixel reading $(0.40, 0.30, 0.10)$.
\item\label{ex:sortbug:3} Bands listed $\mathtt{wl} = (600, 800, 700)$ (unsorted, not descending). Give
  $\mathtt{np.argsort(wl)}$, the sorted list, and the matrix $\mat{M}$ of the fixed \code{resample} for the
  grid $(650, 750)$. Apply it to the pixel $(1, 3, 2)$ in listed order.
\item\label{ex:sortbug:4} Through the old generic resample v0.2.0 caught 15.8\% of the 909 real objects at the top
  1\%, through the product path 44.2\%. If bench-v2.1's FT02c (35.1\%) and F001 (33.1\%) are compared with
  $p = 0.34$ over 2,817 objects, how many standard errors apart are they? Is that difference larger or smaller
  than the bug's effect on v0.2.0's real-bench catch, in standard errors of the real bench ($p = 0.30$, $n = 909$)?
\item\label{ex:sortbug:5} The repo's resample tests passed for weeks with the bug in place. Explain why, and write
  the one-line assertion you would add inside \code{resample} to make the old code fail loudly instead.
\end{checkyourself}
\endgroup

\nextchapter{ch:ftruns}
